Una persona común, al ver un código QR, no piensa nada malo sobre él. En primer lugar, porque no tiene forma de leerlo ni de interpretarlo a simple vista; en segundo lugar, porque ya se acostumbró a encontrarlo en todas partes. Ocurre algo parecido a lo que pasa con un cuchillo en una cocina: de tanto usarlo para cocinar, se olvida que también puede ser un arma. El código QR se percibe como una herramienta de acceso rápido al menú de un restaurante, al pago de un parqueo o al formulario de una institución, y casi nunca como un posible vector de ataque.

Esa confianza tiene un costo. El código QR fue diseñado en 1994 para rastrear componentes en la industria automotriz y está estandarizado en la norma ISO/IEC~18004 \cite{iso18004_2015,iso18004_2024}. Su adopción masiva convirtió un mecanismo industrial en una puerta de entrada cotidiana a contenido digital, presente en comercios, transporte, plataformas de gobierno y sistemas de pago \cite{krombholz2014qr}. Sin embargo, el proceso de escaneo no muestra al usuario el destino del enlace antes de abrirlo. A diferencia de un hipervínculo en pantalla, que puede inspeccionarse antes de tocarlo, el contenido de un QR permanece oculto hasta que el teléfono ya lo está procesando, y la mayoría de los sistemas operativos móviles abren la URL sin aplicar ningún análisis de seguridad \cite{krombholz2014qr,vidas2013qrishing}.

\section{Definición del problema}

Los actores maliciosos aprovechan esta dinámica mediante el \emph{quishing}, término que combina QR y phishing. La técnica consiste en distribuir códigos QR manipulados que llevan a sitios fraudulentos, a la descarga de software malicioso o a formularios que capturan credenciales \cite{insikt2024qrcode,unit42_2023qr}. Frente al phishing tradicional tiene tres ventajas para el atacante. El enlace no es visible antes del escaneo, por lo que el usuario pierde las señales de alerta que le daría una URL sospechosa. Los filtros de correo electrónico, que analizan texto pero no el contenido de una imagen, no lo detectan \cite{unit42_2023qr}. Y puede desplegarse en el mundo físico con una simple calcomanía pegada encima de un código legítimo.

El problema dejó de ser marginal. Insikt Group documentó un crecimiento cercano al 248~\% en las menciones de quishing en fuentes de inteligencia entre el tercer y el cuarto trimestre de 2023, y registró su integración en plataformas de \emph{phishing-as-a-service} como Tycoon~2FA y Greatness \cite{insikt2024qrcode}. Que el ataque se venda como servicio indica que se profesionalizó y que ya no requiere conocimiento técnico para ejecutarlo.

Las deficiencias del ecosistema actual pueden identificarse en dos niveles. A nivel del usuario, existe una percepción generalizada de los códigos QR como elementos neutros o seguros, lo que reduce la vigilancia ante una posible manipulación \cite{krombholz2014qr}. A nivel tecnológico, los lectores integrados en los sistemas operativos no incorporan ningún análisis del contenido antes de redirigir, de modo que toda la carga de detectar el engaño recae en la persona que escanea.

Existen aplicaciones que intentan cubrir este vacío. Herramientas como \emph{Safe QR Scan} y \emph{SafeQR Reader}, disponibles en las tiendas de aplicaciones, verifican la URL antes de abrirla, pero su mecanismo central consiste en contrastar el enlace contra listas negras de dominios conocidos. Este enfoque tiene una limitación documentada desde hace más de una década: las listas negras cubren una fracción pequeña de los sitios de phishing en las primeras horas de vida del ataque, justo cuando ocurre la mayor parte del daño \cite{sheng2009empirical}. En el contexto del quishing, donde el atacante puede registrar un dominio nuevo para cada campaña, depender solo de lo que ya fue reportado deja pasar precisamente los ataques más recientes.

Entonces, el problema que aborda este trabajo puede plantearse así: un usuario de teléfono inteligente no tiene forma de saber a dónde lo lleva un código QR antes de que su dispositivo lo abra, y las soluciones disponibles dependen casi por completo de listas que no conocen las amenazas nuevas.

\section{La propuesta: QRisk}

QRisk es un prototipo de aplicación móvil que se coloca en medio de ese intervalo. Lee el código QR, muestra a dónde apunta en realidad, lo analiza mediante una cascada de filtros y devuelve un veredicto en lenguaje sencillo antes de que el usuario decida si abre el sitio. La regla que guía todo el diseño es que la aplicación nunca visita la URL por su cuenta: hasta que la persona lo pida, el enlace es solo texto en pantalla.

El sistema se desarrolló entre dos personas, cada una a cargo de un módulo. El módulo de servidor, a cargo de Fabiola Contreras, implementa el motor de análisis en la nube con los filtros heurístico (F2.1), de comportamiento (F2.2) y de vista previa visual (F2.3). El módulo que documenta este trabajo, a cargo de la autora, comprende la aplicación móvil completa: el escaneo y la clasificación del contenido, la integración del filtro de reputación F1, el flujo de presentación del veredicto, el historial local, el reporte de códigos maliciosos y la tarjeta compartible de advertencia, conocida en el proyecto como \emph{Story Card}.

\section{Organización del documento}

Continuando con esta introducción, el capítulo~\ref{chap:objetivos} presenta los objetivos aprobados en el protocolo, y el capítulo de justificación explica por qué vale la pena intervenir en el momento del escaneo. El marco teórico reúne los conceptos de seguridad, de desarrollo móvil y de usabilidad que sustentan las decisiones del prototipo, mientras que los antecedentes revisan los trabajos previos y delimitan el vacío que el proyecto busca cubrir. El capítulo de alcance fija lo que el trabajo incluye y lo que deja fuera.

Los capítulos centrales describen la metodología seguida, el desarrollo del prototipo con sus decisiones técnicas y los resultados obtenidos. El documento cierra con las conclusiones, las recomendaciones para trabajo futuro y los anexos con la evidencia de las mediciones.
